\section{Related Work}
\label{sec:related}

\para{Self-critique and self-reflection in LLMs.}
Chain-of-thought prompting~\citep{wei2022chain} demonstrated that intermediate reasoning steps improve performance on complex tasks. \citet{shinn2023reflexion} extended this to iterative self-reflection, showing that language agents can learn from verbal self-evaluations without weight updates. \citet{lightman2023lets} showed that process-level supervision (per-step verification) outperforms outcome-level supervision for mathematical reasoning. \citet{asai2023self} trained models to self-critique via special reflection tokens integrated into generation. \citet{turpin2024language} demonstrated that chain-of-thought explanations can be unfaithful to models' actual reasoning processes, motivating the need to study internal representations rather than surface-level outputs. Unlike these works, which study self-critique at the behavioral level, we directly measure its impact on internal representations.

\para{Representation geometry and probing.}
\citet{marks2024geometry} established that truth and falsehood are linearly represented in LLM hidden states, with consistent directions across topics that support causal interventions. \citet{zou2023representation} introduced representation engineering, using contrastive stimuli and PCA to extract and control concept directions (honesty, truthfulness, emotion) in the residual stream. \citet{polo2026emergent} discovered that concept manifolds undergo transient geometric restructuring during reasoning, distinguishing between information \emph{retention} (high probe accuracy) and computational \emph{readiness} (high manifold capacity). The ``Truth as a Trajectory'' framework~\citep{tat2026} showed that reasoning validity can be detected from layer-wise displacement trajectories rather than static activation snapshots. We draw on these methods---particularly cosine similarity analysis, linear probing, and CKA---to characterize how self-critique reshapes the geometry of internal representations.

\para{Self-reflection vectors and activation steering.}
Most closely related to our work, \citet{zhu2025emergence} demonstrated that self-reflection is encoded as a separable direction in hidden states across multiple model families. They extract ``self-reflection vectors'' via contrastive difference-in-means and show these vectors are domain-invariant (cosine similarity $>$0.95 across tasks). Enhancement via activation steering yields up to +12 percentage point accuracy gains. While their work establishes that reflection \emph{has} a geometric signature, they do not track how representations change through the critique process or compare against matched control text. We complement their findings by quantifying the \emph{magnitude}, \emph{directionality}, and \emph{layer specificity} of critique-induced drift relative to a paraphrase baseline.
